% Section 27.7: The current picture | Part: 6 | Depends on: sec:stokes, sec:ga-calculus | Status: done | Original lines: 3898-3923
%!TEX root = ../../main.tex
\subsection{The current picture}\label{sec:rescurrent}

\begin{itemize}
\item \emph{Riemann surfaces.} If $\eta$ is a meromorphic $1$-form on a compact Riemann surface $S$ (no boundary), Stokes with small discs removed gives $\sum_{p\in S}\Res_p\eta=0$: the sum of all residues vanishes. Combined with Poincar\'e/Serre duality this is a key ingredient in the proofs of the Riemann--Roch theorem and in the theory of abelian differentials. Note that this statement does not even need a metric: $\eta$ is a $1$-form and $\Res_p\eta$ is intrinsic, which is why residues are defined for forms $f\,\dd z$ and not for functions $f$ (the residue of the function $f$ changes under a change of coordinate $z\mapsto h(z)$, that of the form $f\,\dd z$ does not).
\item \emph{Several complex variables.} The kernel $1/(z-z_0)$ is replaced by the Bochner--Martinelli kernel, and residues by Poincar\'e--Leray residues of closed forms with poles along hypersurfaces; both are proved with Stokes' theorem and are expressed through de Rham cohomology of the complement of the singular set.
\item \emph{Currents.} Formulae such as $\partial_{\bar z}(1/\pi z)=\delta_0$ and the Poincar\'e--Lelong formula $\frac i\pi\partial\bar\partial\log|f|=[\{f=0\}]$ treat the residue as the current of integration over the singular set (Section~\ref{sec:current}, item~3), exactly as Stokes' theorem is treated in geometric measure theory \cite{GriffithsHarris}.
\end{itemize}

\begin{center}
\small
\renewcommand{\arraystretch}{1.3}
\begin{tabular}{@{}p{4.6cm}p{9.6cm}@{}}
\toprule
\textbf{Complex analysis} & \textbf{Differential forms}\\
\midrule
$f$ holomorphic & $f\,\dd z$ closed; $\dd f$ is $\C$-linear\\
Cauchy's theorem & Stokes for a closed form on $\Omega$\\
Cauchy--Riemann equations & $\operatorname{div}=0$ and $\operatorname{curl}=0$ (real/imaginary parts of $w\,\dd z$)\\
Residue $\Res_af$ & period of $f\,\dd z$ over a small circle $=$ de Rham pairing\\
Residue theorem & Stokes on $\Omega$ minus discs; boundary $=$ outer curve $-$ small circles\\
Winding number, argument principle & $\frac1{2\pi}\oint f\pb\dd\theta=$ degree of a map to $S^1$\\
Cauchy--Pompeiu, $\partial_{\bar z}(1/\pi z)=\delta$ & Stokes with $\dd\omega\neq0$; Gauss law in the plane\\
\bottomrule
\end{tabular}
\end{center}
